% !TEX root = ../main.tex

%----------------------------------------------------------------------------------------
% CHAPTER 3 - METODOLOGIA I PLANIFICACIÓ
%----------------------------------------------------------------------------------------

\chapter{Metodologia i planificació}
\label{Ch:Methodology}

%----------------------------------------------------------------------------------------

\section{Metodologia de desenvolupament}

S'ha escollit una \textbf{metodologia àgil iterativa} amb les següents característiques:

\subsection{Enfocament incremental}

El desenvolupament s'ha estructurat en iteracions que van incrementant funcionalitat:

\begin{enumerate}
    \item \textbf{Fase 1 - Investigació i configuració}: Investigació de les APIs de Socrata, verificació de disponibilitat de dades, configuració de l'entorn de desenvolupament.
    
    \item \textbf{Fase 2 - Canonada ETL base}: Implementació dels extractors de dades, transformadors de dades i loaders de fitxers.
    
    \item \textbf{Fase 3 - Frontend - Dashboard 1 \& 2}: Creació del servei de dades, hooks personalitzats, i els primers dos panells de visualització (mapa i evolució temporal).
    
    \item \textbf{Fase 4 - Frontend - Dashboard 3 \& 4}: Implementació dels panells de correlació clima-aigua i matriu d'alertes de sequera.
    
    \item \textbf{Fase 5 - Automatització i desplegament}: Configuració de GitHub Actions per a execució diària de l'ETL i construcció del frontend.
\end{enumerate}

\subsection{Justificació de l'enfocament àgil}

\begin{itemize}
    \item \textbf{Flexibilitat}: Permet adaptar-se si les APIs canvien o si sorgeixen nous requisits.
    \item \textbf{Entrega contínua}: Cada iteració proporciona valor incrementat (una canonada ETL funcional, un panell de visualització, etc.).
    \item \textbf{Feedback ràpid}: Es pot provar cada component de forma independent.
\end{itemize}

%----------------------------------------------------------------------------------------

\section{Pla de treball i cronograma}

El pla de treball s'ha estructurat en cinc fases principals, basant-se en una metodologia àgil iterativa que permet rebre feedback continuat i adaptar-se a canvis de requisits o limitacions tècniques descobertes durant el desenvolupament. Cada fase inclou objectius clars, lliurables tangibles i criteris d'acceptació verificables.

\begin{itemize}
    \item \textbf{Setmanes 1-2}: Investigació de dades, configuració del repositori, estudi de les APIs.
    \item \textbf{Setmanes 3-5}: Desenvolupament de la canonada ETL (extractors, transformadors, loaders).
    \item \textbf{Setmanes 6-8}: Desenvolupament del frontend (dashboards 1-2, hooks, serveis).
    \item \textbf{Setmanes 9-10}: Desenvolupament dels dashboards 3-4, integració de navegació.
    \item \textbf{Setmanes 11-12}: Configuració de GitHub Actions, documentació, testing.
\end{itemize}

Els elements destacables del cronograma són:
\begin{itemize}
    \item La Fase 2 (ETL) i Fase 3 (Frontend base) van executar-se en paral·lel durant les setmanes 4-5, aprofitant la independència tècnica dels dos subsistemes.
    \item La Fase 4 (Frontend avançat) va beneficiar-se dels components i serveis creats en la Fase 3, reduint el temps necessari per a la implementació del Dashboard 3 i 4.
    \item Les setmanes 11-12 van dedicar-se predominantement a tasques de qualitat (proves, documentació) en lloc de desenvolupament de noves funcionalitats, assegurant l'estabilitat del lliurament final.
    \item S'hi van assignar 12 hores de tutories (10\%) per a revisions de progrés i ajusts del pla de treball, essencials en una metodologia àgil.
\end{itemize}

%----------------------------------------------------------------------------------------

\section{Tasques i estimació temporal}

Les tasques planificades per a cada fase, amb el temps estimat basat en l'experiència prèvia amb projectes similars i el criteri de complexitat tècnica, són les següents:

\subsection{Fase 1: Investigació i configuració inicial}
\begin{itemize}
    \item Anàlisi detallada de les APIs de Socrata de la Generalitat de Catalunya: identificació d'endpoints rellevants per a embassaments i precipitació, verificació de límits de taxa, examen de formats de resposta (JSON) - 8 hores
    \item Estudi de la documentació tècnica disponible i domini de proves per a validar la qualitat i completitud de les dades - 6 hores
    \item Configuració del repositori Git, estructura de carpetes del projecte i definició de les convencions de codi - 4 hores
    \item Instal·lació i configuració de l'entorn de desenvolupament (Node.js 18, npm 9, Git) - 3 hores
    \item Creació del pla de prova inicial per a validar la connexió bàsica amb les APIs - 3 hores
    \item Redacció de l'anàlisi de viabilitat inicial - 6 hores
    \item \textbf{Total Fase 1:} 30 hores
\end{itemize}

\subsection{Fase 2: Desenvolupament de la canonada ETL}
\begin{itemize}
    \item Implementació del mòdul d'extracció: funcions per a fer peticions HTTP asincròniques a les APIs de Socrata amb maneig d'errors i reintents - 12 hores
    \item Desenvolupament del mòdul de transformació: normalització de dades (conversió d'unitats, format de dates, maneig de valors nul·ls) - 15 hores
    \item Implementació del mòdul de càrrega: emmagatzematge eficient de dades processades en fitxers JSON estructurats per a consum frontend - 10 hores
    \item Creació de l'orquestrador ETL que coordina l'execució seqüencial dels tres mòduls - 5 hores
    \item Implementació del sistema de logging detallat per a monitorització i depuració - 4 hores
    \item Proves unitàries de cada mòdul ETL amb Jest (90\% de cobertura de codi) - 8 hores
    \item Integració i proves de tota la canonada ETL amb simulació de condicions de xarxa adverses - 6 hores
    \item \textbf{Total Fase 2:} 60 hores
\end{itemize}

\subsection{Fase 3: Desenvolupament del frontend - Dashboards 1 i 2}
\begin{itemize}
    \item Configuració del projecte React amb Vite 5, React 18 i React Router 6 - 4 hores
    \item Implementació del servei de consum de dades: wrapper amb axios per a carregar fitxers JSON locals amb maneig d'estats de càrrega i error - 6 hores
    \item Creació de hooks personalitzats per a gestió d'estat global (Zustand) - 5 hores
    \item Desenvolupament del Dashboard 1 - Mapa interactiu: integració de Leaflet 1.9, markers d'embassaments, controls de filtratge, KPIs dinàmics - 18 hores
    \item Desenvolupament del Dashboard 2 - Evolució temporal: Recharts, selector de dates, comparació múltiple, optimització per a grans conjunts de dades - 20 hores
    \item Estilització amb SCSS seguint metodologia BEM - 8 hores
    \item Proves d'integració frontend amb mock de dades JSON - 6 hores
    \item \textbf{Total Fase 3:} 67 hores
\end{itemize}

\subsection{Fase 4: Desenvolupament del frontend - Dashboards 3 i 4}
\begin{itemize}
    \item Desenvolupament del Dashboard 3 - Correlació clima-aigua: gràfic de dispersió, algoritme d'agregació, coeficient de Pearson, controls interactius - 16 hores
    \item Desenvolupament del Dashboard 4 - Matriu d'alertes de sequera: sistema semàfor, taula responsive, algorisme de tendència, classificació automàtica - 14 hores
    \item Implementació del sistema de navegació entre dashboards amb manteniment d'estat - 5 hores
    \item Optimització de rendiment: memoització de components costosos, càrrega diferida - 8 hores
    \item Proves d'usabilitat amb 5 usuaris externs (tècnics i no tècnics) - 6 hores
    \item \textbf{Total Fase 4:} 49 hores
\end{itemize}

\subsection{Fase 5: Automatització, desplegament i documentació}
\begin{itemize}
    \item Configuració de GitHub Actions per a workflow de CI/CD: ETL programat (cron: 0 2 * * *), desplegament frontend, notificacions - 8 hores
    \item Preparació del medi de desplegament: variables d'entorn, optimització d'emmagatzematge, cache-control - 6 hores
    \item Redacció de la memòria del TFG segons les normes de la EPSG - 20 hores
    \item Creació de la documentació tècnica d'ús i manteniment del sistema - 8 hores
    \item Elaboració del vídeo demostratiu (4 minuts 30 segons) mostrant funcionalitats clau - 4 hores
    \item Proves finals d'integració i validació de requisits en entorn de producció simulat - 6 hores
    \item \textbf{Total Fase 5:} 52 hores
\end{itemize}

\subsection{Resum de l'estimació temporal}
\begin{table}[h]
\centering
\caption{Estimació de temps per fase}
\begin{tabular}{|l|c|r|}
\hline
\textbf{Fase} & \textbf{Tasques principals} & \textbf{Hores estimades} \\
\hline
1 & Investigació i configuració inicial & 30 \\
\hline
2 & Desenvolupament de la canonada ETL & 60 \\
\hline
3 & Frontend - Dashboards 1 i 2 & 67 \\
\hline
4 & Frontend - Dashboards 3 i 4 & 49 \\
\hline
5 & Automatització, desplegament i documentació & 52 \\
\hline
\textbf{TOTAL} & & \textbf{258} \\
\hline
\end{tabular}
\label{tab:time-estimation}
\end{table}

L'estimació total de 258 hores s'acosta molt a la inversió real de 250 hores documentada durant el desenvolupament, indicant una bona capacitat de predicció de l'esforç necessari.

%----------------------------------------------------------------------------------------

\section{Paquets de treball}

El desenvolupament s'ha dividit en tres paquets de treball (PT) principals, cadascun corresponent a un subsistema funcional del projecte:

\textbf{PT1: Canonada de Processament de Dades (ETL)}
\begin{itemize}
    \item \textbf{Objectiu}: Construir un pipeline fiable i escalable per a l'extracció, transformació i càrrega de dades hidrològiques des de les fonts públiques de la Generalitat de Catalunya.
    \item \textbf{Lliurables}: Scripts d'extracció parametritzats, funcions de transformació verificades, mòdul d'emmagatzematge optimitzat, orquestrador ETL amb maneig d'errors robust.
    \item \textbf{S'explica als}: Capítols 7 (decisions ETL), 9 (implementació del pipeline ETL).
\end{itemize}

\textbf{PT2: Aplicació Web de Visualització}
\begin{itemize}
    \item \textbf{Objectiu}: Desenvolupar una interfície web intuïtiva i responsiva que permeti l'exploració interactiva de dades hidrològiques mitjançant diversos paradigmes de visualització.
    \item \textbf{Lliurables}: Aplicació React de pàgina única amb quatre dashboards interconnectats, servei de consum de dades optimitzat, sistemes de gestió d'estat i navegació.
    \item \textbf{S'explica als}: Capítols 7 (decisions frontend), 9 (implementació del frontend), 10 (validació de requisits funcionals).
\end{itemize}

\textbf{PT3: Infraestructura d'Automatització i Desplegament}
\begin{itemize}
    \item \textbf{Objectiu}: Implementar un sistema d'integració i desplegament continu que asseguri la frescor de les dades i la disponibilitat del servei amb mínima intervenció manual.
    \item \textbf{Lliurables}: Workflows de GitHub Actions configurats, scripts de desplegament parametritzats, mecanismes de monitorització i alertes.
    \item \textbf{S'explica als}: Capítols 7 (decisions de desplegament), 10 (desplegament del sistema, validació de requisits no funcionals).
\end{itemize}

Aquesta descomposició en paquets de treball va facilitar la distribució de les responsabilitats i permetre un seguiment més granular del progrés durant el desenvolupament.

%----------------------------------------------------------------------------------------

\section{Eines utilitzades}

\subsection{Control de versions}
\begin{itemize}
    \item \textbf{Git}: Sistema de control de versions distribuït.
    \item \textbf{GitHub}: Repositori remot i plataforma de col·laboració.
    \item \textbf{GitHub Actions}: CI/CD per a l'automatització de l'ETL.
\end{itemize}

\subsection{Desenvolupament backend (ETL)}
\begin{itemize}
    \item \textbf{Node.js 18}: Runtime de JavaScript.
    \item \textbf{npm}: Gestor de paquets.
    \item \textbf{axios}: Client HTTP per a fer peticions a les APIs.
    \item \textbf{dotenv}: Gestió de variables d'entorn.
\end{itemize}

\subsection{Desenvolupament frontend}
\begin{itemize}
    \item \textbf{React 18}: Framework frontend.
    \item \textbf{Vite}: Bundler ràpid i modern.
    \item \textbf{React Router}: Navegació entre pàgines.
    \item \textbf{Recharts}: Biblioteca de gràfics.
    \item \textbf{Leaflet}: Biblioteca de mapes interactius.
    \item \textbf{Sass}: Preprocessador CSS.
\end{itemize}

\subsection{Documentació i testing}
\begin{itemize}
    \item \textbf{LaTeX}: Preparació de la memòria del projecte.
    \item \textbf{Visual Studio Code}: Editor de codi.
    \item \textbf{Chrome DevTools}: Debugging del frontend.
\end{itemize}

%----------------------------------------------------------------------------------------

\section{Ús d'eines d'intel·ligència artificial}

\subsection{Assistència de codificació}

S'ha fet ús de \textbf{GitHub Copilot} com a assistent de codificació, principalment per a:

\begin{itemize}
    \item \textbf{Boilerplate de React}: Generació d'estructures de components, hooks, i renders.
    \item \textbf{Estils SCSS}: Creació de fulls d'estil responsives.
    \item \textbf{Funcions de transformació de dades}: Lògica de filtrat i agrupament de dades.
\end{itemize}

\subsection{Declaració d'ús de la IA}

\begin{itemize}
    \item \textbf{Nivell d'ús}: Aproximadament 30-40\% del codi a nivell de suggeriments (propostes completades o parcials).
    \item \textbf{Verificació}: Totes les suggerències van ser revisades, verificades i probades per al context específic del projecte.
    \item \textbf{Lògica original}: Els algoritmes de processament de dades, arquitectura del sistema i decisions de disseny són originals del desenvolupador.
    \item \textbf{Documentació}: La majoria de comments i documentació de codi va ser generada pel desenvolupador, amb assistència puntual de la IA per a claredat.
\end{itemize}

\textbf{Nota}: L'ús de la IA va ser \textbf{sempre transparent} i \textbf{sempre supervisat}. La IA va ser una eina d'assistència per a accelerar la codificació repetitiva, no per a substituir el pensament creatiu o la presa de decisions tècniques del desenvolupador.

%----------------------------------------------------------------------------------------

\section{Desviacions i incidències}

\subsection{Desviacions respecte a la planificació}

\begin{enumerate}
    \item \textbf{Canvi d'enfocament en la visualització de correlació}: Inicialment es va planificar utilitzar mapes de calor per a mostrar la correlació entre precipitació i nivells d'embassament. Durant les proves inicials es va observar que aquesta representació enfosquia les tendències temporals. Es va canviar per a un gràfic de dispersió amb línies de tendència i càlcul del coeficient de correlació, proporcionant una mesura quantitativa més rigorosa. Aquest canvi va requerir 4 hores addicionals però va millorar notablement la qualitat analítica del Dashboard 3.
    
    \item \textbf{Incorporació de mecanisme de predicció tendencial}: El pla original no incloïa cap funcionalitat predictiva dins del Dashboard 4. Es va implementar un algorisme senzill de regressió lineal sobre les últimes 6 mesures per a detectar canvis significatifs, afegint una fletxa de tendència al costat de cada indicador. Aquesta millora va requerir 3 hores addicionals i va ser ben rebuda durant les proves d'usabilitat.
    
    \item \textbf{Inclusió de més estacions meteorològiques}: Originalment es planificava de mostrar només algunes estacions, però les APIs proporcionen 200+ estacions. S'ha inclòs tot el dataset per a màxima cobertura geogràfica.
    
    \item \textbf{Quatre dashboards en lloc de tres}: La viabilitat tècnica va demostrar que era fàcil afegir un quart dashboard (alertes), proporcionant més valor a l'usuari.
\end{enumerate}

\subsection{Incidències i resolucions}

\begin{enumerate}
    \item \textbf{Incidència}: Mida dels fitxers JSON generats és gran (30MB per a embassaments, 70MB per a precipitació).
    \textbf{Resolució}: Els fitxers es serveixen des de \texttt{/public/data} i es carreguen al frontend de forma reactiva (es pot implementar compressió o paginació en futures millores).
    
    \item \textbf{Incidència}: Terser no estava instal·lat al desenvolupament inicial del frontend.
    \textbf{Resolució}: S'instal·la de forma manual com a dependència de desenvolupament.
    
    \item \textbf{Incidència}: Problemes amb caràcters especials (Àlex) en els camins de Windows.
    \textbf{Resolució}: Ús consistent de comandes PowerShell amb \texttt{-LiteralPath} per a evitar problemes de codificació.
\end{enumerate}

%----------------------------------------------------------------------------------------

\section{Conclusions de la planificació}

El pla de treball va resultar eficaç per a guiar el desenvolupament del projecte des de la concepció fins al lliurament, demostrant les següents característiques positives:
\begin{itemize}
    \item La descomposició en fases clares amb objectius verificables va facilitar el seguiment del progrés i la detecció primerenca de desviacions.
    \item L'estimació temporal va resultar prou precisa (error de +3,2\%), indicant una bona comprensió de la complexitat tècnica involucrada.
    \item La flexibilitat inherent a la metodologia àgil va permetre incorporar millores significatives sense comprometre els límits temporals.
    \item La documentació del pla de treball va servir com a eina de comunicació eficient amb el tutor i altres interessats.
\end{itemize}

Les principals aprenentatges negatius relacionats amb la planificació van ser:
\begin{itemize}
    \item \textbf{Subestimació de la complexitat de l'optimització de rendiment frontend}: Tot i que es va preveure temps per a tasques d'optimització, la magnitud de l'esforç necessari per a mantenir una experiència fluida amb conjunts de dades >50K punts va ser major a l'esperat, requerint una redistribució de 5 hores des de tasques de funcionalitats noves a tasques d'optimització durant les darreres setmanes.
    \item \textbf{Dependència excessiva de proves manuals per a validació d'usabilitat}: La planificació inicial va donar massa pes a les proves automatitzades i poc a les proves amb usuaris reals. Durant la Fase 4 es va detectar problemes d'usabilitat que no van ser capturats per les proves unitàries, necessitant una reassignació de 4 hores per a proves d'usabilitat estructurades amb el públic objectiu. En futures planificacions s'haurà de destinar un mínim del 20\% del temps de proves a validació amb usuaris reals des de les fases inicials.
\end{itemize}

En general, el pla de treball va proporcionar un marc sòlid que va permetre adaptar-se a les circumstàncies canviants mantenint l'objectiu final: crear una plataforma útil i accessible per a la monitorització de recursos hídrics a Catalunya.

%----------------------------------------------------------------------------------------
